\documentclass[11pt]{article}
\usepackage[margin=1.1in]{geometry}
\usepackage{amsmath,amssymb,amsthm}
\usepackage{hyperref}
\newtheorem{theorem}{Theorem}[section]
\theoremstyle{definition}
\newtheorem{conjecture}[theorem]{Conjecture}
\theoremstyle{remark}
\newtheorem*{remark}{Remark}
\title{A Machine-Verified Proof of TLMC Conjecture \#590:\\ Theta Series of Numerical Semigroups}
\author{Batch 11 solution submission}
\date{September 12, 2026}
\begin{document}
\maketitle
\begin{abstract}
We \emph{prove} a conjecture of the Last Math Competition: for a numerical semigroup $S$ with Frobenius number $F$, the theta series $\Theta_S(q) = \sum_{s \in S} q^s$ satisfies $\Theta_S(q) - \frac{1}{1-q} = -\sum_{g \in \mathrm{Gaps}} q^g$, a polynomial with at most $F$ monomials (the gaps all lie in $\{1, \dots, F\}$ since $0 \in S$ and $n > F \Rightarrow n \in S$). The Lean~4 formalization defines the gap polynomial as an explicit monomial sum, realizes $1/(1-q)$ via the geometric power series (\texttt{ones}, with the kernel-verified identity $(1 - q)\cdot\texttt{ones} = 1$), proves the coefficient identity of $\Theta_S + G = 1/(1-q)$, and bounds the support. Compilation is clean; standard axioms only; no \texttt{sorry}.
\end{abstract}
\section{The conjecture}
\begin{conjecture}[TLMC-00000000590]\label{conj:590}
For a numerical semigroup $S$ with Frobenius number $F$: $\Theta_S(q) - \frac{1}{1-q}$ is a polynomial with at most Frobenius-many monomials (a sparsity bound).
\end{conjecture}
\section{The proof}
\begin{theorem}\label{thm:590}
Conjecture~\ref{conj:590} is \textbf{true}. More precisely, for $S \subseteq \mathbb{N}$ with $0 \in S$, $F \notin S$ and the conductor property $n > F \Rightarrow n \in S$:
\[
\Theta_S(q) + G(q) = \frac{1}{1-q}, \qquad G(q) := \sum_{\substack{g \le F \\ g \notin S}} q^g,
\]
and $G$ has at most $F$ monomials.
\end{theorem}
\begin{proof}
The identity is coefficientwise: at $n$, the $n$-th coefficient of $\frac{1}{1-q}$ is $1 = [n \in S] + [n \le F \wedge n \notin S]$ (if $n \in S$ the first indicator is $1$ and the second is $0$; if $n \notin S$ then $n \le F$ by the conductor property, and the second indicator is $1$). For the sparsity: $\mathrm{Supp}(G) = \{g \le F : g \notin S\}$ has at most $F$ elements, since it is contained in $\{1, \dots, F\}$ ($0 \in S$ excludes $0$) — formally, the gap set is a proper subset of $\{0, \dots, F\}$ by the exclusion of $0$, hence has at most $F$ elements.
\end{proof}
\section{Machine verification}
\texttt{Batch11.lean} (namespace \texttt{TLMCBatch11}) contains: \texttt{ones} and \texttt{one\_sub\_mul\_ones} (the geometric series with $(1 - q)\cdot\texttt{ones} = 1$, proved via antidiagonal coefficient sums), \texttt{thetaS}, \texttt{gapPoly} with \texttt{coeff\_gapPoly} (its coefficients are the gap indicators, via \texttt{sum\_eq\_single\_of\_mem}), and the main theorem \texttt{tlm590} packaging the identity $\Theta_S + \iota(G) = \texttt{ones}$ and the bound $\#\mathrm{Supp}(G) \le F$. Compilation via \texttt{lake env lean Batch11.lean} exits $0$ with zero errors and warnings; \texttt{\#print axioms} reports at most \texttt{[propext, Classical.choice, Quot.sound]}; no \texttt{sorry}.
\end{document}
